\ifja
\chapter{現代の欺瞞と主権の回復}
\label{chap:introduction}

\section{設計された認知的不協和}
ソフトウェア工学や物理、数学に携わる技術者が、いわゆる「簿記」や「税務」の教科書を初めて開いたときに覚える激しい違和感と拒絶反応は、決して彼らの知性の欠如に起因するものではない。むしろ逆である。第一原理（First Principles）からの演繹、論理的一貫性、そして保存則の普遍性を信じる訓練を受けた健全な頭脳であるからこそ、そこに蔓延する「設計された不整合」にスタックオーバーフローを起こしているにすぎない。

古くは中世イタリアの商人に起源を持つとされる複式簿記は、本来、代数的な保存則そのものであった。しかし現代の資本主義社会においてそれは、資格制度、士業ギルド、そして高額なサブスクリプションを課すクラウドSaaSベンダーによって幾重もの「認知的トラップ」で包囲され、神聖化されたブラックボックスへと改悪されてしまった。

彼らは、技術者が論理的に世界を理解できないよう、意図的に次のような認知的拷問を仕掛けている。

\begin{enumerate}
    \item \textbf{対称性の意図的な破壊（「借方・貸方」の亡霊）：}\\
    工学や物理において、ノードに対する流入と流出は単なる符号付きスカラー量（流入を $+I$、流出を $-I$）として定式化される。明治期に福澤諭吉らが「Debit / Credit」を「借方 / 貸方」と翻訳して以来、日本の教育現場では「借りたわけではないのに借方」「資産の増加は左だが負債の増加は右」という言語的・数学的に破綻した呪文の丸暗記を初心者に強要し続けている。これは概念の本質（符号保存則）を隠蔽し、学習者を迷宮に閉じ込めるための第一の障壁である。

    \item \textbf{状態空間の不要な散乱（名詞による目眩まし）：}\\
    10万円未満なら消耗品費、10万円以上30万円未満なら一括償却または少額減価償却資産、30万円以上なら固定資産台帳へ登録して耐用年数テーブルを参照……。これらは数学的には単なる金額と日付を入力とする「閾値判定関数（Step Function）」にすぎない。しかし業界はこれに夥しい名詞のラベルを貼り付け、例外パッチを継ぎ接ぎすることで、「専門家にしか扱えない高等技術」であるかのように演出している。

    \item \textbf{キャッシュフローの神秘化（間接法の煙幕）：}\\
    「利益が出ているのに現金がない（黒字倒産）」「営業キャッシュフローの計算には減価償却費を足し戻し、売掛金の増減を差し引き……」といった実務教育は、あたかもキャッシュフロー計算書（C/F）が既存の財務諸表とは独立した別個の魔術であるかのように喧伝する。しかしその実態は、資産空間における現金基底への単なる射影と移項代入にすぎない。手続きをわざわざ間接法という複雑な調整ステップで包囲し、本質的なスループットの把握を困難にさせている。

    \item \textbf{恐怖の煽動と小作農化（Rent-seeking）：}\\
    「少しでも間違えれば税務署に差し押さえられる」「法律は複雑怪奇だから素人が触るな」という恐怖を刷り込むことで、個人の主権を放棄させる。結果として技術者は、銀行口座の認証情報を怪しいクラウド家計簿に差し出し、重厚なWeb GUIの前で毎月小作料（月額サブスクリプション）を払い、画面上のボタンをポチポチと押すだけのデジタル小作農へと転落させられる。
\end{enumerate}

\section{第一原理への還元：キルヒホッフ則と方程式}
この茶番劇の皮を剥ぎ取った後に残る真実は、高校の理数科目レベル（一次方程式と数列）の道具立てで完全に記述できる。

\subsection{キルヒホッフの電流則（KCL）と同型の複式簿記}
あらゆる会計取引の本質は、有向グラフにおけるノード（勘定科目）間のフロー保存則、すなわち\textbf{キルヒホッフの電流則（Kirchhoff's Current Law: KCL）}と同型である。

任意の会計取引グラフにおけるノード集合を $V$、取引 $e$ におけるノード $v \in V$ への資金流入量を $I_e(v)$ と定義するとき、閉鎖系における総代数和は常にゼロとなる。

\begin{equation}
    \sum_{v \in V} I_e(v) = 0 \quad (\forall e)
\end{equation}


勘定科目の集合を頂点、資金の移動を有向枝とする閉鎖系において、出た分だけ入り、全体の代数和は常に厳密にゼロとなる。出納における「借方（左）」とは単なる $+I$（資金の運用・投下先）であり、「貸方（右）」とは単なる $-I$（資金の源泉・調達元）にすぎない。

\subsection{キャッシュフロースループットの代数的演繹}
キャッシュフロー計算書（C/F）という独立した怪物は存在しない。資産（$\text{Assets}$）を現金（$\text{Cash}$）と非現金資産（$\text{NonCash}$：売掛金、在庫、機材、不動産等）に直和分解するだけで、キャッシュの真の運動方程式は完全に演繹される。

\begin{equation}
    \text{Assets} = \text{Cash} + \text{NonCash} \implies \text{Cash} \equiv \text{Assets} - \text{NonCash}
\end{equation}
\begin{equation}
    \text{Assets} = \text{Profit} + \text{Liabilities} + \text{Equity}_0
\end{equation}
\begin{equation}
    \mathbf{\text{Cash} = \text{Profit} + \text{Liabilities} + \text{Equity}_0 - \text{NonCash}}
\end{equation}
\begin{equation}
    \mathbf{\Delta \text{Cash} = \text{Profit} + \Delta \text{Liabilities} + \Delta \text{Equity} - \Delta \text{NonCash}}
\end{equation}


「黒字倒産」の正体は、$\text{Profit} > 0$ であっても、非現金資産の膨張分（$\Delta \text{NonCash}$）が利益を上回った瞬間に $\Delta \text{Cash} < 0$ へ沈むという、単純な代数符号の逆転にすぎない。

\else
\chapter{The Modern Deception and Restoration of Sovereignty}
\label{chap:introduction}

\section{Engineered Cognitive Dissonance}
The severe cognitive dissonance experienced by software engineers, physicists, and mathematicians when first opening traditional accounting textbooks is not due to any intellectual deficit. On the contrary, minds trained in deductive reasoning from First Principles, logical consistency, and universal conservation laws suffer an immediate stack overflow when confronted with deliberately engineered incoherence.

Double-entry bookkeeping, originating with medieval merchants, was purely an algebraic conservation law. Modern commercial institutions, professional guilds, and subscription-gated SaaS vendors have surrounded it with cognitive traps, degrading it into a mystified black box.

\begin{enumerate}
    \item \textbf{Deliberate Destruction of Symmetry (The Ghost of "Debit and Credit"):}\\
    In physics and engineering, inflows and outflows are signed scalars ($+I$ for inflow, $-I$ for outflow). Dogmatic rules forcing beginners to memorize that "debit means left, credit means right" obscure the underlying algebraic sign-conservation law.
    \item \textbf{State-Space Fragmentation (Obfuscation via Terminology):}\\
    Depreciation tiers and threshold classifications are mathematically nothing more than Step Functions over price and temporal domains.
    \item \textbf{Mystification of Cash Flow (The Indirect Method Smokescreen):}\\
    The standard indirect cash flow statement merely obfuscates a straightforward vector projection onto the liquid cash basis.
    \item \textbf{Rent-Seeking via Manufactured Fear:}\\
    Cultivating fear of regulatory non-compliance induces individuals to forfeit operational autonomy and become digital tenant farmers paying monthly rents to web GUIs.
\end{enumerate}

\section{Reduction to First Principles: Kirchhoff's Current Law}
Every financial transaction is strictly isomorphic to a conservation flow across a directed graph, formally equivalent to \textbf{Kirchhoff's Current Law (KCL)}.

Let $V$ denote the set of account nodes, and $I_e(v)$ denote the algebraic flow entering node $v \in V$ for transaction event $e$. In any closed system, the net divergence is identically zero:
\begin{equation}
    \sum_{v \in V} I_e(v) = 0 \quad (\forall e)
\end{equation}


"Debit" is strictly $+I$ (application/sink of capital), while "Credit" is strictly $-I$ (source of capital).

\subsection{Algebraic Derivation of Cash Throughput}
The Statement of Cash Flows is not an independent construct. By decomposing total Assets into Cash and Non-Cash, the dynamic throughput equation is derived algebraically:
\begin{equation}
    \text{Assets} = \text{Cash} + \text{NonCash} \implies \text{Cash} \equiv \text{Assets} - \text{NonCash}
\end{equation}
\begin{equation}
    \text{Assets} = \text{Profit} + \text{Liabilities} + \text{Equity}_0
\end{equation}
\begin{equation}
    \mathbf{\text{Cash} = \text{Profit} + \text{Liabilities} + \text{Equity}_0 - \text{NonCash}}
\end{equation}
\begin{equation}
    \mathbf{\Delta \text{Cash} = \text{Profit} + \Delta \text{Liabilities} + \Delta \text{Equity} - \Delta \text{NonCash}}
\end{equation}


Insolvent growth (profitable bankruptcy) is simply the sign reversal of $\Delta \text{Cash} < 0$ when non-cash asset accumulation ($\Delta \text{NonCash}$) exceeds net operational throughput.
\fi
